\prefacesection{Abstract}

In non-stationary environments, time series anomaly detection is challenging because normal behavior changes can resemble true anomalies. Fixed detectors or a single global threshold may misclassify concept drift as abnormal behavior, while overly adaptive detectors may incorporate true abnormal data into the normal model. This report introduces the AnDri system, an interactive web platform for jointly detecting anomalies and drift in univariate time series. AnDri and the baseline methods provide anomaly scores for each time point by comparing local behavior with learned normal patterns. The system then converts these scores into final labels through threshold judgment and user feedback.

This project focuses on the implementation of the AnDri system and its interactive workflow. The constructed platform supports data set upload, training area selection, anomaly labeling, candidate detection result review, detector execution, normal mode checking, and result comparison. The platform uses Python Flask backend, combined with JavaScript and Apache ECharts for visualization, and integrates AnDri with baseline detectors such as NormA, SAND, and DAMP. Through a human-machine collaborative threshold optimization process, using partial user annotations, Bayesian Gaussian mixture model fitting, and active candidate review to adjust anomaly thresholds; for AnDri, different normal mode groups can be independently optimized for thresholds.

This system has been applied to different time series datasets to demonstrate how the platform supports anomaly review, drift exploration, threshold optimization, and comparison with baseline methods. These case studies show that the implemented interface can present time series anomaly detection algorithms' outputs, user feedback, normal mode information, and evaluation metrics in a single workflow. The comparison results also show that AnDri is more effective than baseline methods in concept drift scenarios, as it can better distinguish normal changes caused by drift from true abnormal events. The system not only reports automatic detection results but also helps users analyze why suspicious areas are determined as anomalies or abnormal behavior changes. The report also records the system architecture, visualization design, integration with baseline methods, and deployment experience on school servers.